\begingroup
\ArithmeticHead{306}{L}{ANALYSIS. NACHLASS.}
\NumberedItem{15}
\par\indent Die Siebentheilung führt auf folgende Gleichung
\[\begin{aligned}
b&=\left(\frac{1-2x+2x^4-{.}{.}{.}}{1+2x+2x^4+{.}{.}{.}}\right)^2\\[6pt]
A&=\left(\frac{1+2x^7+2x^{28}+{.}{.}{.}}{1+2x+2x^4+{.}{.}{.}}\right)^2\\[6pt]
B&=\left(\frac{1-2x^7+2x^{28}-{.}\quad}{1+2x+2x^4+{.}\,{.}\;}\right)^2
\end{aligned}\]
\clearpage
\ArithmeticHead{307}{R}{{\scriptsize ZUR THEORIE DER TRANSSCENDENTEN FUNCTIONEN. I.}}
\SourceLeaf{314}
\[\begin{aligned}
&A^8-\frac47A^6+\frac{16-32bb}{49}A^5-\frac{30}{343}A^4+\frac{32-64bb}{2401}A^3\\
&\quad-\frac{20+768bb-768b^4}{16807}A^2+\frac{48-2144bb+6144b^4-4096b^6}{823\,543}A\\
&\quad-\frac1{823543}=0
\end{aligned}\]
\[\begin{aligned}
&B^8-\frac47bbB^6+\frac{16b^3-32b}{49}B^5-\frac{30b^4}{343}B^4+\frac{32b^5-64b^3}{2401}B^3\\
&\quad-\frac{20b^6+768b^4-768b^2}{16807}B^2+\frac{48b^4-2144b^5+6144b^3-4096b}{823\,543}[B]\\
&\quad-\frac1{823543}=0.
\end{aligned}\]
\noindent Wenn man in der ersten Gleichung statt $bb$, $1-bb$ und statt $A$, $-A$ setzt,
so wird sie nicht geändert.
\clearpage
\endgroup
